\section{РАЗРАБОТКА МОДУЛЯ ИНТЕГРАЦИИ И РАЗВЁРТЫВАНИЕ ОКРУЖЕНИЯ}

В разделе описаны разработанный модуль интеграции формата Vortex в Apache Paimon,
функциональное тестирование артефакта в контейнерах, развёртывание окружения
нагрузочного тестирования на кластере средствами Ansible, методика проведённых
экспериментов и проблемы, выявленные и решённые в ходе работы.

\subsection{Модуль интеграции формата Vortex в Apache Paimon}

Модуль реализован как набор подмодулей Maven в составе Apache Paimon (интеграция
выполнена в ответвлении репозитория, на основе версии Apache Paimon 1.4):
\texttt{paimon-vortex-jni} --- обёртки нативных методов и загрузчик библиотеки;
\texttt{paimon-vortex-format} --- реализация интерфейса \texttt{FileFormat} и
связанных классов; нативная библиотека Vortex собирается из исходных текстов на
Rust и упаковывается в артефакт. Логически модуль состоит из четырёх частей: путь
записи, путь чтения, конвертер предикатов и слой доступа к нативной библиотеке.

\textbf{Слой доступа к нативной библиотеке.} Класс \texttt{NativeFileMethods}
объявляет \texttt{native}-методы: открытие файла по URI с параметрами хранилища
(\texttt{open}), получение числа строк и схемы (\texttt{rowCount}, \texttt{dtype}),
запуск сканирования с проекцией столбцов, сериализованным предикатом и диапазоном
строк (\texttt{scan}), закрытие (\texttt{close}), а также операции уровня каталога
(перечисление и удаление файлов). Загрузчик \texttt{NativeLoader} извлекает
библиотеку под текущую платформу из ресурсов артефакта и подгружает её. Передача
больших массивов данных через границу JNI выполняется без копирования: нативная
сторона экспортирует декодированный пакет в раскладке Apache Arrow через Arrow C
Data Interface, а JVM-сторона импортирует его в \texttt{VectorSchemaRoot}.

\textbf{Путь записи.} \texttt{VortexFileFormat} создаёт \texttt{FormatWriterFactory},
которая по пути в хранилище порождает писатель: строки Paimon (\texttt{InternalRow})
преобразуются в пакеты Arrow заданного размера (\texttt{write.batch-size}) и
передаются нативному писателю Vortex; по завершении файла формируется объект
статистик по столбцам (минимум, максимум, число null) для записи в манифест Paimon.
Размер пакета записи существенен: он же определяет размер пакета, возвращаемого при
чтении одним нативным вызовом, поэтому от него зависит число пересечений границы JNI
при последующем сканировании. В работе используется увеличенный относительно
умолчания размер пакета, что снижает накладные расходы JNI при чтении (см.
подраздел~\ref{subsec:issues}).

\textbf{Путь чтения.} \texttt{VortexFileFormat} создаёт \texttt{FormatReaderFactory};
её метод порождает \texttt{VortexRecordsReader} по пути файла, проекции столбцов,
(опционально) предикату и параметрам хранилища. Читатель открывает нативный файл,
запускает сканирование (передавая список проецируемых столбцов, сериализованный
предикат и --- при наличии --- список номеров строк для выборочного чтения) и
получает итератор нативных массивов. Метод \texttt{readBatch} извлекает очередной
нативный массив, экспортирует его в Arrow \texttt{VectorSchemaRoot} (с
переиспользованием ранее выделенной структуры), преобразует средствами
\texttt{ArrowBatchReader} в итератор \texttt{InternalRow} и возвращает движку
вместе с информацией о позиции строки в файле (необходима Paimon для
deletion-векторов и слияния). Нативные ресурсы (массив, аллокатор Arrow, итератор,
дескриптор файла) освобождаются при закрытии читателя или при освобождении пакета.

\textbf{Конвертер предикатов.} \texttt{VortexPredicateConverter} преобразует
предикат Paimon в нативное выражение Vortex. Поддержаны операции сравнения
(\(=, \ne, <, \le, >, \ge\)), проверка на null, конъюнкция и дизъюнкция для всех
типов данных, для которых это имеет смысл (логический, целочисленные, строковый,
десятичный, дата, временная метка). Непереводимые предикаты не проталкиваются ---
фильтрация по ним остаётся на стороне движка; корректность при этом не нарушается.
Выражение сериализуется и передаётся нативной библиотеке, которая применяет его при
сканировании до пересечения границы JNI, что сокращает объём данных, передаваемых в
JVM, на запросах с селективным предикатом.

\textbf{Многопоточность.} Параллелизм обеспечивается на уровне Paimon и Flink: чтение
таблицы разбивается на разделённые задачи (splits) --- по бакетам и файлам, ---
каждая обрабатывается своим читателем в отдельном потоке (subtask) вычислительного
движка; запись в бакет однопоточна, параллелизм записи равен числу бакетов. Внутри
одного читателя нативное сканирование Vortex также может использовать несколько
потоков. Объекты модуля, разделяемые между задачами (фабрики), не имеют изменяемого
состояния; объекты, привязанные к одной задаче (читатель, аллокатор Arrow,
нативные дескрипторы), потоконебезопасны и используются строго из одного потока, что
соответствует контракту интерфейсов Paimon. Параметры памяти настраиваются с учётом
того, что нативная сторона использует прямые (off-heap) буферы вне кучи JVM ---
этому посвящён отдельный пункт подраздела~\ref{subsec:issues}.

Последовательность взаимодействия компонентов при чтении таблицы Paimon с форматом
Vortex показана на рисунке~\ref{fig:seq-read}.

\begin{figure}
\centering
\resizebox{\textwidth}{!}{%
\begin{tikzpicture}[
    font=\footnotesize,
    >={Stealth[length=2mm]},
    head/.style={draw, rounded corners=1.5pt, align=center, minimum height=8mm, minimum width=20mm, fill=black!4},
    msg/.style={->, thick},
    ret/.style={->, dashed},
  ]
  % --- координаты времени жизни ---
  \def\xA{0}      % Flink
  \def\xB{3.6}    % Paimon FileFormat
  \def\xC{7.2}    % VortexFileFormat
  \def\xD{10.8}   % VortexRecordsReader
  \def\xE{14.4}   % Vortex (Rust) via JNI
  \def\ytop{0}
  \def\ybot{-10.6}
  % --- головы участников ---
  \node[head] (A) at (\xA,\ytop) {Flink\\(движок)};
  \node[head] (B) at (\xB,\ytop) {Apache Paimon\\\texttt{FileFormat}};
  \node[head] (C) at (\xC,\ytop) {\texttt{VortexFile}\\\texttt{Format}};
  \node[head] (D) at (\xD,\ytop) {\texttt{VortexRecords}\\\texttt{Reader}};
  \node[head] (E) at (\xE,\ytop) {Vortex (Rust)\\через JNI};
  % --- линии жизни ---
  \foreach \x in {\xA,\xB,\xC,\xD,\xE} {
    \draw[dashed, black!55] (\x,-0.45) -- (\x,\ybot);
  }
  % --- сообщения ---
  % 1
  \draw[msg] (\xA,-1.2) -- node[above]{\texttt{createReaderFactory(проекция, предикат)}} (\xB,-1.2);
  % 2
  \draw[msg] (\xB,-1.9) -- node[above]{\texttt{createReaderFactory(\dots)}} (\xC,-1.9);
  % 3 self: convert predicate
  \draw[msg] (\xC,-2.6) -- ++(1.0,0) -- ++(0,-0.5) -- node[right, xshift=2pt]{\texttt{convert()} — предикат $\to$ нативное выражение} ++(-1.0,0);
  % 4 return factory
  \draw[ret] (\xC,-3.7) -- node[above]{\texttt{FormatReaderFactory}} (\xB,-3.7);
  % 5 new reader
  \draw[msg] (\xB,-4.4) -- node[above]{\texttt{new VortexRecordsReader(путь, тип, предикат)}} (\xD,-4.4);
  % 6 open + scan
  \draw[msg] (\xD,-5.1) -- node[above]{\texttt{open(uri)}; \texttt{scan(столбцы, предикат)}} (\xE,-5.1);
  \draw[ret] (\xE,-5.7) -- node[above]{итератор нативных массивов} (\xD,-5.7);
  % --- loop frame ---
  \draw[black!60] (-1.5,-6.2) rectangle (15.9,-9.3);
  \node[draw, fill=white, font=\scriptsize\bfseries, anchor=north west] at (-1.5,-6.2) {loop [пока есть пакеты]};
  % loop messages
  \draw[msg] (\xA,-6.9) -- node[above]{\texttt{readBatch()}} (\xD,-6.9);
  \draw[msg] (\xD,-7.6) -- node[above]{\texttt{next()}; \texttt{exportToArrow()}} (\xE,-7.6);
  \draw[ret] (\xE,-8.2) -- node[above]{Arrow \texttt{VectorSchemaRoot} (без копирования)} (\xD,-8.2);
  \draw[ret] (\xD,-8.9) -- node[above]{пакет строк (\texttt{InternalRow})} (\xA,-8.9);
  % --- close ---
  \draw[msg] (\xA,-9.8) -- node[above]{\texttt{close()}} (\xD,-9.8);
  \draw[msg] (\xD,-10.4) -- node[above]{освобождение нативных ресурсов} (\xE,-10.4);
\end{tikzpicture}%
}
\caption{Диаграмма последовательности: чтение таблицы Apache Paimon с файловым форматом Vortex}
\label{fig:seq-read}
\end{figure}
